%% kapa: exam=1 round=19 year=2008 period=1 form=A subject=appraisal_law subject_ko=부동산관계법규 seq=24 number=64
%% source: Q-Net 감정평가사 기출문제 (cst00309) artlSeq=1001962 file=extracted/1008/1001962/2008년 제19회 감정평가사 1차시험 기출문제/1교시.hwp
%% conversion: kordoc 4.24.51 -> tools/gen1.py (auto); review=proofread-sonnet; flags=none
\begin{question}{64}
부동산 가격공시 및 감정평가에 관한 법률에서 규정하는 과징금에 관한 설명으로 옳은 것은?
\begin{choices}
  \item 국토해양부장관이 과징금을 부과하는 경우에는 위반행위의 내용과 정도, 위반행위의 기간과 위반횟수, 위반행위로 취득한 이익의 규모를 고려하여야 한다.
  \item 국토해양부장관은 과징금납부의무자가 납부기한 내에 과징금을 납부하지 아니한 경우에는 납부기한일부터 납부한 날까지의 기간에 대하여 가산금을 징수할 수 있다.
  \item 지정한 기간 내에 과징금을 납부하지 아니하였을 때에는 민사집행법에 의하여 징수한다.
  \item 과징금의 부과처분에 이의가 있는 자는 그 처분을 통보받은 날로부터 60일 이내에 이의를 신청할 수 있다.
  \item 과징금납부의무자가 과징금 납부기한을 연장받거나 분할납부를 하려는 경우에는 납부기한 14일 전까지 신청하여야 한다.
\end{choices}
\end{question}
\answer{1}
